\subsection{Specialization and the exponent-one algebraic rule}
\label{subsec:ftb-key-specialization-descent}

In this subsection we assume the exact connected-family and local
analytic family inputs used in
Proposition~\ref{prop:ftb-key-complex-equivariance}, together with
the already constructed normalized powered rule \(\rho\).
Thus \(N(g)\ge1\), \(\rho\) is functorial on characteristic-zero
fields, is a unit on every DVR carrying a key datum, and
\[
                             \rho(x)=t(x)^{N(g)}
                                      \quad(x/\mathbb C).
 \tag{S1}\label{eq:ftb-key-powered-normal-value}
\]
The value \(t(x)=\det(M)^4/c^8\) is independent of normal form and
is equivariant under \emph{all} field maps
\(\mathbb C\to\mathbb C\).
These hypotheses retain the outstanding family producers stated
in the preceding subsection; the argument below does not assert
their existence a second time.

\begin{srlemma}[Countable avoidance and a smooth curve through a point]
\label{lem:ftb-countable-smooth-through}
Let \(R\) be a domain finitely generated over \(\mathbb Z\), with
a map \(a:R\to\mathbb C\).
There are a smooth integral finitely generated
\(\mathbb C\)-algebra \(E\), an injection \(i:R\hookrightarrow E\),
and a complex point \(e:E\to\mathbb C\), with \(ei=a\).
One can take \(\dim E\le1\).
Moreover, if \(E_1\) is any integral finite-type
\(\mathbb C\)-algebra and \(T\subset E_1\setminus\{0\}\) is
countable, some complex point of \(E_1\) is nonzero on all of \(T\).
\end{srlemma}
\begin{proof}
First prove the last assertion.
Noether normalization gives an integral finite inclusion
\(P=\mathbb C[X_1,\ldots,X_d]\hookrightarrow E_1\).
For each nonzero \(r\in T\), integrality and the domain
hypothesis give \(0\ne b_r\in P\cap(r)\).
For example a monic relation for \(r\) of least degree has
nonzero constant term, which belongs to \((r)\).
The nonvanishing set of a nonzero polynomial on \(\mathbb C^d\)
is open and dense: vanishing on an open polydisc would make
the polynomial zero.
The Baire theorem gives \(c\in\mathbb C^d\) at which every
\(b_r\) is nonzero.
A maximal ideal of \(E_1\) above the evaluation ideal at \(c\)
has residue field \(\mathbb C\), by integrality and algebraic
closedness.  Its evaluation cannot kill any \(r\).

To construct the smooth curve, start with
\(\mathbb C\otimes_{\mathbb Z}R\) and the evaluation induced by \(a\).
Flat going down supplies a prime below its evaluation kernel
whose contraction to \(R\) is zero.
Its quotient \(B\) is an integral finite-type
\(\mathbb C\)-algebra, the map \(R\to B\) is injective, and the
given evaluation descends.
It remains to choose a curve through this particular point
without losing that injection.

Choose a finite integral Noether normalization
\(P=\mathbb C[X_1,\ldots,X_d]\hookrightarrow B\).
For every nonzero image of an element of \(R\), choose a
nonzero \(b_r\in P\cap(r)\) as above.
There are only countably many such \(r\).
Let \(c_0\in\mathbb C^d\) be the image of the prescribed point,
and choose \(c_1\) at which all \(b_r\) are nonzero.
Use the line substitution
\[
                        X_j\longmapsto (c_0)_j+
                                      ((c_1)_j-(c_0)_j)T.
 \tag{S2}\label{eq:ftb-line-substitution}
\]
Its kernel is prime and is contained in the contraction of
the prescribed maximal ideal.
Going down for the integral extension from the normal
polynomial ring supplies a prime \(Q\) of \(B\), below that
maximal ideal, lying over this line kernel.
It avoids every \(r\): otherwise \(b_r\) would restrict
identically to zero on the line, contradicting its value
at \(c_1\).
The quotient \(B/Q\) has dimension at most one.
Indeed its contraction quotient is either \(\mathbb C\)
or \(\mathbb C[T]\), and integral incomparability prevents
a longer prime chain upstairs.

Take the integral closure of \(B/Q\) in its fraction field.
It is finite, since \(B/Q\) is a finite-type curve over a field.
It is a Noetherian normal domain of dimension at most one,
hence regular, and is smooth over the perfect field
\(\mathbb C\).
The given point lifts by lying over, with residue field
\(\mathbb C\).
Its coordinate ring is \(E\); the avoidance of all nonzero
\(r\) gives the asserted injection \(R\hookrightarrow E\).
This is the source's line construction and normalization;
it is not an arbitrary smoothing of the original base.
\end{proof}

\begin{srlemma}[Finite-coefficient and countable local-PID descent]
\label{lem:ftb-key-countable-local-pid-specialization}
Let \(V\) be a characteristic-zero local PID domain.
A key datum over \(V\), together with a countable set of its
coefficients, descends to a countable local PID subring
\(V_0\subset V\).
The inclusion is a local homomorphism and is faithfully flat.
If only a finite set of extra coefficients is prescribed and
\(V\) is Noetherian local, the key datum descends to a domain
finitely generated over \(\mathbb Z\), contained in \(V\);
the latter subring need not be local.
\end{srlemma}
\begin{proof}
We first record the ring construction.
For a countable \(T\subset V\), let \(k_0\) be the subfield
of \(\operatorname{Frac}V\) generated by \(T\), and put
\(V_0=V\cap k_0\).
This is countable.
Divisibility of its elements is reflected from \(V\):
if \(b=ac\) in \(V\) and \(a\ne0\), then \(c=b/a\in k_0\),
so \(c\in V_0\).
In particular the inclusion reflects units, and \(V_0\)
is local.
It is a PID.  For an ideal \(I\subset V_0\), choose among
the ideals \(aV\), \(a\in I\), a maximal one, using
Noetherianity of \(V\).
Total divisibility in the local PID \(V\) shows that its
chosen \(a\) divides every \(b\in I\) in \(V\).
Divisibility reflection gives \(I=aV_0\).
The \(V_0\)-module \(V\) is torsion free, hence flat over
this PID, and a flat local inclusion is faithfully flat.

Here is the geometric descent used with this construction.
Choose a very ample projective embedding of the abelian
scheme, from a sufficiently high polarization power.
On the finitely many standard affine charts its defining
ideals have finite sets of polynomial generators.
Adjoin their coefficients.
The closed image in projective space over the resulting
subring pulls back to exactly the original scheme:
on each chart the ideal is the extension of those generators.
Adjoin the finitely many coefficients needed for the
identity, multiplication and inverse maps, and for the
line's finite transition functions and their inverses.
These maps are descended on finite affine refinements;
their polynomial identities are checked after the
faithfully flat inclusion.
Thus the group laws descend.
Flatness, properness, smoothness, the relative dimension,
the line isomorphisms giving symmetry and rigidification,
relative ampleness, and the principal polarization
condition descend by faithful flatness.
This gives the actual polarized abelian scheme over a
countable local PID of the displayed kind.

Choose a key frame on that local base.
Pull it back to \(V\).
The prescribed key datum differs from it by its line
isomorphism, a theta scalar \(a\in V^\times\), and a Lie
matrix \(M\in\operatorname{GL}_g(V)\).
Adjoin the scalar and matrix entries to the countable
coefficient set and repeat the intersection construction.
Unit reflection gives \(a^{-1}\) and
\(\det(M)^{-1}\) in the new local subring.
Change the descended frames by these actual parameters.
The line, its rigidification, the theta section and all
ordered tangent vectors now pull back to the original
key datum, as asserted.

For the finite-type assertion, first descend the polarized
projective data to finitely many coefficients.
Near the contraction of the closed point choose a
nonzero theta section and a Lie basis.
Invert a single element not vanishing there.
Its image is a unit of \(V\); adjoin its inverse, the
required theta scalar and inverse, the Lie matrix entries
and inverse determinant, and the specified finite
coefficient set.
The same actual rescaling gives a key datum over this
finite-type subring.
This is the finite-coefficient construction in the source,
not a claim that an arbitrary frame extends before the
necessary coefficients and inverses are adjoined.
\end{proof}

\begin{srproposition}[Specialization of the unpowered complex value]
\label{prop:ftb-key-normal-value-specializes}
Let \(V\) be a local PID domain with a key datum \(y/V\).
Let \(f,f':V\to\mathbb C\), with \(f\) injective.
If \(u\in V\) and \(t(f^*y)=f(u)\), then
\[
                                t((f')^*y)=f'(u).
 \tag{S3}\label{eq:ftb-key-in-families}
\]
\end{srproposition}
\begin{proof}
Descend \(y,u\) by the finite-type assertion of
Lemma~\ref{lem:ftb-key-countable-local-pid-specialization}, to a finite-type
domain \(R\subset V\) and a key datum \(y_R\).
Apply Lemma~\ref{lem:ftb-countable-smooth-through} to
\(f'|_R\), obtaining \(i:R\hookrightarrow E\) and \(e:E\to\mathbb C\).
Write \(z=i^*y_R\).
The smooth-base powered-unit lemma gives
\(b\in E^\times\) whose evaluation at every complex point is
\(t(z_s)^{N(g)}\).

If \(s:E\to\mathbb C\) does not kill any nonzero \(i(r)\),
then \(si:R\to\mathbb C\) is injective.
The two injections \(f|_R\) and \(si\) admit embeddings
\(\alpha,\beta:\mathbb C\to\mathbb C\) into a common
complex field copy with
\(\alpha f|_R=\beta si\).
This follows by extending the two embeddings of the
countable fraction field, using transcendence bases and
algebraic closures.
Equivariance of \(t\) under these possibly discontinuous
field maps, and its independence of normal form, give
\[
                              t(z_s)=s(i(u)).
\]
Indeed the two key pullbacks become identical after the
common extension, and \(\beta\) is injective.

It follows that \(i(u)^{N(g)}=b\) in \(E\).
If their difference were nonzero, countable avoidance
would produce a complex point not killing that
difference or any nonzero \(i(r)\), contradicting the
last equality and the defining equality for \(b\).
Hence \(i(u)\) is a unit.
Choose one of these generic complex points \(s_0\).
The root-fixed-by-one-point lemma, now on the connected
smooth integral complex curve \(\operatorname{Spec}E\),
compares \(t(z_s)\) with \(s(i(u))\) at every point:
their \(N(g)\)-th powers agree, and their ratio is fixed
by its value \(1\) at \(s_0\).
At \(e\) this gives \(t(z_e)=e(i(u))=f'(u)\).
The fibre \(z_e\) is the prescribed special key pullback,
so this is \eqref{eq:ftb-key-in-families}.
\end{proof}

\begin{srtheorem}[The exponent-one field rule and its DVR units]
\label{thm:ftb-key-rule-exponent-one}
Under the inputs of this subsection there is a unique
rule \(\kappa_K(x)\in K^\times\), for every characteristic-zero
field \(K\) and key datum \(x/K\), which agrees with
\(t(x)\) over \(\mathbb C\), is functorial under all field
maps and all key pullbacks, and has
\[
                   a^8\kappa(a\theta,Mv)=\det(M)^4\kappa(\theta,v).
 \tag{S4}\label{eq:ftb-key-rule-exponent-one-homog}
\]
For every DVR \(V\), fraction field \(K\) of characteristic
zero, and key datum \(y/V\), its generic value is the
image of a unit \(u\in V^\times\).
For every characteristic-zero residue-field quotient
\(r:V\twoheadrightarrow k\),
\[
                                 \kappa_k(r^*y)=r(u).
 \tag{S5}\label{eq:ftb-key-rule-exponent-one-unit}
\]
For a normalized theta function and the ordered tangent
frame \((2\pi i)^{-1}e_j\), the value is \(1\).
\end{srtheorem}
\begin{proof}
We first define the value over a countable field \(k\).
Embed \(k\) in \(\mathbb C\), and pull the key datum to \(x\).
If \(\alpha\) is any automorphism of \(\mathbb C\) fixing \(k\),
apply \eqref{eq:ftb-key-in-families} with the local PID
\(\mathbb C\), \(f=1\), and \(f'=\alpha\).
It gives the transformed value \(\alpha(t(x))\).
Because \(\alpha\) fixes \(k\), cancellation of their common
\(k\)-pullbacks gives an actual key-pullback isomorphism over
\(1_{\mathbb C}\) between \(\alpha^*x\) and \(x\).
Apply the same assertion with \(f'=1\) using that isomorphism;
it gives the value \(t(x)\).
Thus \(\alpha(t(x))=t(x)\).
The fixed field of all such automorphisms is \(k\).
For an element algebraic over \(k\), a distinct conjugate
extends to an automorphism unless the element lies in \(k\),
since the characteristic is zero.
For a transcendental element, changing a transcendence
basis element and extending to algebraic closures gives
a nonfixing automorphism.
Therefore \(t(x)=s(u)\) for a unique \(u\in k\).
Equation \eqref{eq:ftb-key-in-families}, with \(V=k\), makes
this same \(u\) work at every embedding of \(k\).

For an arbitrary \(K\), descend the key datum to a
countable subfield \(k\subset K\), using
Lemma~\ref{lem:ftb-key-countable-local-pid-specialization} and then taking
the fraction field.
Define its value by the preceding \(u\).
This is independent of descent.
Indeed put two countable models in one countable
intermediate field, adjoin the finitely many coefficients
of their common abelian-variety isomorphism and line
isomorphism, and descend these isomorphisms.
Their equality with the given rigidification, section
and tangent vectors descends by faithful flatness.
Embed this common countable field in \(\mathbb C\).
Both proposed values are then the same complex normal
value; injectivity proves equality before embedding.
This also proves functoriality, by pushing one countable
model through a field map.
The complex normal value is nonzero, so the value belongs
to \(K^\times\).
The same descent including \(a,M\) and their inverses
transfers the complex normal-form identity
\eqref{eq:ftb-known-frame-change} to \(K\), giving
\eqref{eq:ftb-key-rule-exponent-one-homog}.

The relation with the powered rule is
\[
                             \kappa_K(x)^{N(g)}=\rho_K(x).
 \tag{S6}\label{eq:ftb-key-root-powered-algebraic}
\]
Check it over a countable model embedded in \(\mathbb C\),
using \eqref{eq:ftb-key-powered-normal-value}, and then
use functoriality.
Over the DVR \(V\), the right side is a unit \(b\in V^\times\).
The element \(\kappa_K(x)\) is integral over \(V\), being
a root of \(T^{N(g)}-b\).
Normality of \(V\) puts it in \(V\); its positive power
being a unit makes it a unit \(u\).

We must also identify the specialization of this chosen
root, not just its power.
Descend \(y,u\) to a countable local PID subring
\(V_0\subset V\), with local faithfully flat inclusion.
The generic countable fraction field embeds in \(\mathbb C\).
By the definition and uniqueness of its value, the
generic complex value is the image of \(u\).
Equation \eqref{eq:ftb-key-in-families} now identifies the
value under every map \(V_0\to\mathbb C\), including maps
factoring through a characteristic-zero residue field.
It therefore represents the value on that residue
pullback in the countable definition.
Pushing to \(k\) proves
\eqref{eq:ftb-key-rule-exponent-one-unit}.
For a normalized theta datum, \(M=I,c=1\), so its value
is \(1\).
All uniqueness assertions follow by the same common
countable model argument.
\end{proof}

\paragraph{Exact scope and source concordance.}
This theorem gives \(\mathrm{Homog}_1\), \(\mathrm{Funct}\),
\(\mathrm{UnitSp}\), and \(\mathrm{ThetaOne}\).
In \(\mathrm{UnitSp}\) the DVR may have positive residue
characteristic; its unit assertion still holds, while the
stated value comparison is only for characteristic-zero
residue fields, exactly as the field rule requires.
The separate archimedean norm identity is not inferred
from the algebraic root argument.

The prefix below is
\srcid{src/Development/AbelianVarietyTheory/Proofs/ThetaLawProofs/}.
The countable avoidance, chosen line and smooth normalization
are \srcid{ExistsAlgHomNeZeroOfCountable.lean}
(\srcid{9dc863ab22573a559ae8c249f33d533a7860cf49}),
\srcid{ExistsDominantComponent.lean}
(\srcid{8e6078832242f5fd87c9969273e88a9b1a4402df}),
\srcid{ExistsPrimeDimensionLEOneAvoiding.lean}
(\srcid{f56c1b74e6815aeea22584eaf146e6252e747ef8}),
and \srcid{ExistsSmoothNormalisation.lean}
(\srcid{2e88aff00b5fb1cfda8566ac7c21761fdce68dff}).
The exact local subring is
\srcid{ExistsCountableLocalPidSubring.lean}
(\srcid{80f4e931782aab0999790894204dc2b851005194}).
Finite projective descent on standard affine charts is
\srcid{VeryAmpleByDescendsCountableSubring.lean}
(\srcid{880d25ecad16a9516141cd62e4c9feee229e7b53}).
Its actual key-data assembly is
\srcid{KeyDataDescendCountableSubrings.lean}
(\srcid{d06ec6cf161d65aff27cc78f9f3693252a1b58d2}),
using finite-coefficient descent for morphisms, line bundles,
group laws and polarizations.
Those are ordinary projective finite-presentation and fpqc
descent operations with the hypotheses stated in the proof;
no uniform bound for their coefficient fields is needed.

The specialization proof, including the finite-type
rescaling of the \emph{actual} key datum, is
\srcid{ExistsSmoothThrough.lean}
(\srcid{fa1c6cdf3595de268b05d93491d66de7c70a4532}).
The common-countable-isomorphism argument, invariant field
definition, root identity and DVR specialization are
\srcid{ExistsCommonCountable.lean}
(\srcid{f748a7e68705d47f6d415d95bf2e20e2370d2577}).
The final wrapper is
\srcid{../UniformizationProofs/KeyFormulaOfConnFam.lean}
(\srcid{c6229bc642171a481eef723c9bdfffc7b9f68a55}).
Its actual connected-family and analytic-continuity
hypotheses remain precisely those retained above.
